How does a force determine a trajectory? Newton's second law gives an
acceleration, not a position. Two integrations connect force to motion, and
two initial conditions select one trajectory from a family of possible motions.

\section{Model and initial conditions}
Consider a particle of constant mass $m>0$ moving along a fixed $x$ axis in an
inertial reference frame. Let the signed, constant \emph{net} force along that
axis be $F_0$. Other forces must already be included in this sum. For example,
a horizontal applied force may be the net force along a frictionless track;
gravity and the normal force then cancel vertically.

Newton's second law is
\begin{equation}
 F_0 = m\frac{dv}{dt} = m\frac{d^2x}{dt^2},
 \qquad a_0 := \frac{F_0}{m}.
\end{equation}
We specify position and velocity at the same initial time $t_0$:
\begin{equation}
 x(t_0)=x_0, \qquad v(t_0)=v_0.
\end{equation}
The time $t_0$ need not be zero. Write $\tau=t-t_0$ for elapsed time.
These assumptions describe nonrelativistic motion with no drag or changing
mass. Constant acceleration is valid only while the net force stays constant.

\section{First integration: acceleration to velocity}
Integrate $dv/dt=a_0$ with respect to time:
\begin{equation}
 v(t)=a_0t+C_1.
\end{equation}
The derivative of a constant is zero, so differentiation cannot tell us $C_1$.
The initial velocity supplies the missing information:
\begin{equation}
 v_0=a_0t_0+C_1,
 \qquad C_1=v_0-a_0t_0.
\end{equation}
Thus
\begin{equation}
 \boxed{v(t)=v_0+a_0(t-t_0)=v_0+a_0\tau.}
\end{equation}
Equivalently, a definite integral incorporates the condition directly:
\begin{equation}
 v(t)-v_0=\int_{t_0}^{t}a_0\,ds=a_0(t-t_0).
\end{equation}
The integration variable $s$ is a dummy variable; $t$ is the upper limit.

\section{Second integration: velocity to position}
Integrating $dx/dt=a_0t+C_1$ gives a second independent constant:
\begin{equation}
 x(t)=\frac{1}{2}a_0t^2+C_1t+C_2.
\end{equation}
Now use the initial position and the value of $C_1$:
\begin{align}
 x_0 &= \frac{1}{2}a_0t_0^2+(v_0-a_0t_0)t_0+C_2,\\
 C_2 &= x_0-v_0t_0+\frac{1}{2}a_0t_0^2.
\end{align}
Substitution and collection of terms yield
\begin{equation}
 \boxed{x(t)=x_0+v_0(t-t_0)+\frac{1}{2}a_0(t-t_0)^2.}
\end{equation}
With the choice $t_0=0$, the unshifted constants become $C_1=v_0$ and $C_2=x_0$.
In elapsed-time coordinates the constants are always simply the initial values.

The same result follows by a definite integral of the velocity:
\begin{align}
 x(t)-x_0
 &=\int_{t_0}^{t}\bigl[v_0+a_0(s-t_0)\bigr]\,ds\\
 &=v_0\tau+\frac{1}{2}a_0\tau^2.
\end{align}
This exhibits the double integration explicitly:
\begin{equation}
 x(t)=x_0+v_0\tau+
 \int_{t_0}^{t}\left[\int_{t_0}^{s}a_0\,du\right]ds.
\end{equation}
The factor $1/2$ is the integral of the linearly growing velocity increment.
Geometrically it is the triangular area under $v-v_0$ versus elapsed time.

\section{A worked example with a nonzero start time}
Take $m=2\,\mathrm{kg}$, $F_0=4\,\mathrm{N}$, $t_0=2\,\mathrm{s}$,
$x_0=1\,\mathrm{m}$ and $v_0=3\,\mathrm{m\,s^{-1}}$. Then
$a_0=2\,\mathrm{m\,s^{-2}}$, and
\begin{align}
 v(t)&=3\,\mathrm{m\,s^{-1}}+(2\,\mathrm{m\,s^{-2}})\tau,\\
 x(t)&=1\,\mathrm{m}+(3\,\mathrm{m\,s^{-1}})\tau
                   +(1\,\mathrm{m\,s^{-2}})\tau^2.
\end{align}
At clock time $t=5\,\mathrm{s}$, only $\tau=3\,\mathrm{s}$ have elapsed.
The particle is at $x=19\,\mathrm{m}$ with velocity $v=9\,\mathrm{m\,s^{-1}}$.
Its displacement is $18\,\mathrm{m}$, not $19\,\mathrm{m}$.

The constants in the unshifted form are
$C_1=-1\,\mathrm{m\,s^{-1}}$ and $C_2=-1\,\mathrm{m}$.
They encode the chosen clock and coordinate origin; they are not the prescribed
velocity and position at $t_0=2\,\mathrm{s}$.

\section{Three forces, the same initial state}
The accompanying Julia resource holds mass and initial conditions fixed and
compares three constant net forces over $0\leq\tau\leq6\,\mathrm{s}$:
\begin{center}
\begin{tabular}{lrrr}
Case & $F_0\ (\mathrm{N})$ & $a_0\ (\mathrm{m\,s^{-2}})$ & $x(t_0+6\,\mathrm{s})\ (\mathrm{m})$\\
\hline
Speeding up & 4 & 2 & 55\\
Coasting & 0 & 0 & 19\\
Slowing, then reversing & -2 & -1 & 1
\end{tabular}
\end{center}
\begin{center}
\includegraphics[width=0.95\linewidth]{comparison.png}
\end{center}
The zero-force case demonstrates Newton's first law: an object can move while
the net force is zero. A force is required to \emph{change} its velocity.

In the negative-force case, $v=0$ at $\tau=3\,\mathrm{s}$, or
$t=5\,\mathrm{s}$, where $x=5.5\,\mathrm{m}$. Acceleration is still
$-1\,\mathrm{m\,s^{-2}}$ at that instant. The force then drives the particle
backward. This is not a brake that switches off when the particle stops.
After six seconds the displacement is zero, but the distance traveled is
$4.5+4.5=9\,\mathrm{m}$. Signed velocity and speed $|v|$ are different quantities.

The third-law reaction acts on the agent supplying the force, not on the
particle itself; it does not cancel $F_0$ in the particle's force balance.

\section{Checks and a useful time-free relation}
Differentiate the position twice to recover $v(t)$ and $a_0$. At $\tau=0$,
both prescribed initial conditions are recovered. All terms in the position
formula have units of length.

Eliminating elapsed time gives
\begin{equation}
 v^2-v_0^2=2a_0(x-x_0).
\end{equation}
One can check this by substituting the two integrated formulas, without
dividing by $a_0$ or $v$. Multiplying by $m/2$ gives the work--energy relation
\begin{equation}
 \frac{1}{2}mv^2-\frac{1}{2}mv_0^2=F_0(x-x_0).
\end{equation}
Squaring loses the sign of velocity. The time-dependent solution is needed to
distinguish the outward and returning parts of the turning example.

\section{Computational method and further experiments}
The saved results use Julia's standard-library matrix exponential. For
$\mathbf{u}=(x,v,1)^{\mathsf T}$,
\begin{equation}
 \frac{d\mathbf{u}}{dt}=A\mathbf{u},\qquad
 A=\begin{pmatrix}0&1&0\\0&0&a_0\\0&0&0\end{pmatrix},\qquad
 \mathbf{u}(t)=\exp(\tau A)\mathbf{u}(t_0).
\end{equation}
Because $A^3=0$, the exponential terminates at
$I+\tau A+\tau^2 A^2/2$, recovering the two integrations above.
This is evaluation of an exact linear-system solution in floating-point
arithmetic, not a time-stepping approximation. Independent tests compare it
with the integrated formulas, initial conditions, signed area under velocity,
and work--energy balance. The explorer displays 121 saved samples per case,
spaced by $0.05\,\mathrm{s}$; its controls do not run Julia.

\PMlinkexternal{Open the constant-acceleration explorer}{https://images.physicslibrary.org/examples/newton-constant-acceleration/index.html}.
The explorer includes CSV data and a downloadable source project.

Try predicting what changes if the same force acts on twice the mass, if
$v_0=0$, or if the clock origin changes while elapsed time stays fixed.
Edit the local Julia case file to test your predictions. Changing the coordinate
origin changes $x_0$ but neither velocity nor acceleration.
